% =====================================================================
%  Método 8 — Comprimento de dependência e não-projetividade  (estrutura)
%  Origem: células 33, 34 e 35 do caderno
% =====================================================================
\subsection{Comprimento de dependência e não-projetividade}
\label{met:dependencia}

Os métodos anteriores examinavam quais palavras se vinculam a quais. Este examina a distância
que as separa na ordem linear e a eventual sobreposição entre os vínculos. São duas medidas
extraídas da mesma árvore, ambas operando sobre apenas duas informações --- o conjunto de
arcos e a posição de cada palavra ---, sem que intervenham etiquetas, morfologia ou léxico.

A primeira é o comprimento de dependência. Cada palavra $t$, exceto a raiz, vincula-se ao seu
núcleo $\nucleo{t}$; sendo $i(\cdot)$ a posição na frase, o arco apresenta comprimento
$d(t) = \lvert\, i(t) - i(\nucleo{t}) \,\rvert$, e a soma sobre a frase integral fornece
%
\begin{equation}
D(\sent) = \sum_{t \,\neq\, \raiz} d(t)
\qquad\qquad
\bar{D}(\sent) = \frac{D(\sent)}{|A|}
\end{equation}
%
em que $|A|$ designa o número de arcos, excluída a pontuação. Uma vez que $D$ cresce com a
extensão da frase, apenas $\bar{D}$ é comparável entre frases de comprimentos diferentes.

A justificativa da medida é de ordem cognitiva. Um arco vincula duas palavras cuja relação é
necessária à interpretação; ao processar a primeira, o leitor mantém o vínculo em aberto até
encontrar a segunda, e quanto maior o arco, maior o intervalo de retenção. É a hipótese da
minimização do comprimento de dependência~\cite{gibson1998,liu2008}, para a qual há evidência
em larga escala em dezenas de línguas~\cite{futrell2015}.

A segunda medida é a não-projetividade. Tomem-se dois arcos, cada um representado com a
extremidade de menor índice em primeiro lugar, $(i,j)$ e $(k,l)$, com $i<j$ e $k<l$. Os arcos
se cruzam quando exatamente uma extremidade de cada um recai no interior do outro:
%
\begin{equation}
i < k < j < l \qquad\text{ou}\qquad k < i < l < j
\end{equation}
%
Sendo $X(\sent)$ o total de pares que se cruzam, a árvore é projetiva quando $X(\sent) = 0$,
isto é, quando é possível traçar todos os arcos acima da linha sem sobreposição entre eles. A
condição é puramente aritmética sobre índices, não comportando elemento linguístico algum.

O corpus contém duas frases com as mesmas 13 palavras e os mesmos 12 arcos, ambas corretas,
diferindo apenas na ordem. A Tabela~\ref{tab:dependencia} apresenta as duas medidas sobre
elas.

\begin{table}[H]
\centering
\dimtable{
\begin{tabular}{@{}lrrrr@{}}
\toprule
\textbf{Frase} & $D$ & $\bar{D}$ & \textbf{Maior arco} & $X$ \\
\midrule
O vento agravou a seca que intensificou o fogo que destruiu a mata. & 17 & 1,42 & 2 & 0 \\
O fogo que a seca que o vento agravou intensificou destruiu a mata. & 43 & 3,58 & 9 & 0 \\
\bottomrule
\end{tabular}}
\caption{As duas medidas sobre o mesmo conteúdo em ordens distintas.}
\label{tab:dependencia}
\end{table}

Na primeira frase nenhum arco excede 2, pois cada relativa se vincula ao substantivo
imediatamente anterior e a frase se resolve da esquerda para a direita sem retenção. Na
segunda os vínculos se acumulam --- \texttt{fogo} encontra seu verbo nove posições adiante
--- e $\bar{D}$ mais que duplica. Como referência aproximada, a prosa comum em português
situa-se entre 1 e 2; acima de 3, a frase mantém número elevado de vínculos simultaneamente
em aberto.

Este é o aspecto inacessível aos índices da subseção anterior. Ambas as frases encaixam duas
relativas e receberiam a mesma profundidade e o mesmo número de orações, de modo que o que as
distingue é a posição em que as relativas foram inseridas. Já $X$ é nulo em ambas, como
ocorre na quase totalidade dos casos, o que a Seção~\ref{sec:empiricos} retoma.
